% Part 2 chapter carrying "Entanglement, Decoherence, and the Thermodynamic Cost of Classical Records" (Mar 2026, revised Oct 2026), read in full
% 2026-10-02 (PDF text 253 lines; ledger docs/book/read_ledgers/LEDGER_G7_31_Entanglement.md). Corrections applied: PAPER_ERRATA EN1-EN7, GD3, GD4, MP1, MP4.
% Numbers: docs/verification/scripts/verify_entanglement_electroweak.py (E1-E7). Check: docs/verification/quantum/ENTANGLEMENT_CHECK.md.
% Already carried elsewhere and only referenced here: exponent n = 7/2 and its bottom-up running, light writes nothing in flight, tau_IAM at 1 ng (1e-12 kg)
% (Chapter ch:quantumrecords); Bell/CHSH background (ch:nonlocal); where a record becomes irreversible, Bennett (ch:measurement); tau_IAM (ch:gravdec).
\chapter{Entanglement and the cost of records}\label{ch:entanglement}

\section*{The place}
Chapter~\ref{ch:quantumrecords} followed IAM's Law from the quantum-to-classical transition up to the cosmic horizon. This chapter asks what the law
says about entangled systems. Three things are separated: what it leaves unchanged, what it adds, and what an experiment could test. The Bell and
CHSH background is in Chapter~\ref{ch:nonlocal}; where a record becomes irreversible is in Chapter~\ref{ch:measurement}.

\section{Two questions}\label{sec:ent:two}
Two questions are often merged. The first is why entangled correlations are non-classical. The answer is that an entangled pair is described by one
quantum state, not by two systems with pre-existing local properties~\cite{Bell1964}. The second is where, and at what cost, that state gives way to
definite classical outcomes.

Decoherence answers part of the second question. Interaction with an environment suppresses interference and selects pointer
states~\cite{Zurek1981,JoosZeh1985,Zurek2003}. Quantum Darwinism explains why many observers read the same outcome: the record is copied redundantly
into the environment~\cite{Zurek2009}. Neither fixes what the record costs.

IAM's Law does: every irreversible transition from quantum superposition to classical record pays $\kB T\ln2$ per bit to the nearest encoding
surface at the local temperature (Chapter~\ref{ch:iams_law}). In a laboratory the surface is the detector that absorbs the record, and $T$ is its
temperature. Acquiring the record has no minimum dissipation; the $\kB T\ln2$ is paid when the record is erased or the apparatus is
reset~\cite{Landauer1961,Bennett1982}. For bound structure at cosmological scale the nearest shared surface is the horizon, and the cost enters the
horizon entropy (Chapter~\ref{ch:quantumrecords}).

\section{What the law leaves unchanged}\label{sec:ent:unchanged}
The law changes nothing in quantum mechanics and adds no hidden variables. The correlations of a pure shared state are those of standard theory until
an irreversible record is written. For a photon pair in vacuum transit nothing acts in flight (Chapter~\ref{ch:quantumrecords}), so an ideal source
gives $|S|=2\sqrt2$ at any separation. Losses and decoherence in fibre or air are ordinary environmental effects. \derived

In laboratory systems the dominant decoherence is environmental: gas collisions, thermal photons, phonons. The law leaves these channels as they are.
It adds an accounting statement, that every irreversible record carries its Landauer cost at the surface that holds it, and one gravitational channel
for massive superpositions. Separation does not discriminate between the two pictures: the distances reached with photons and with matter qubits are
set by losses, not by sector physics (Chapter~\ref{ch:nonlocal}).

\section{What dephasing does to a Bell test}\label{sec:ent:dephasing}
The gravitational channel of Chapter~\ref{ch:quantumrecords} acts on the which-path superposition of a massive body. It is decoherence in the position, or
pointer, basis: it leaves the populations of the two branches unchanged and multiplies their coherence by a factor $c(t)$ that falls from 1 to 0.
Take a pair entangled in that basis, $(|L\rangle|0\rangle+|R\rangle|1\rangle)/\sqrt2$. With the pointer basis as the $z$ axis, its correlation matrix
is $T={\rm diag}(c,-c,1)$. The largest CHSH value over all settings is $2\sqrt{M}$, with $M$ the sum of the two largest eigenvalues of
$T^{\mathsf T}T$~\cite{Horodecki1995}. Here $M=1+c^2$, so
\begin{equation}
S_{\max}(t)=2\sqrt{1+c(t)^2}.\label{eq:ent:smax}
\end{equation}
\derived{} Two consequences follow. First, $S_{\max}>2$ for every $c>0$: a pair decohering in its pointer basis keeps a Bell violation for as long as
any coherence remains. Second, at $c=0$ the value is exactly 2, not 0. The correlation in the pointer basis survives intact; it has become a classical
correlation between two records. That is what decoherence is expected to leave.

Two other cases fix the reading of a measured $S$. With the settings that are optimal for the pure state, $S=\sqrt2\,(1+c)$, which falls below 2 once
$c<\sqrt2-1=0.414$. \calc{} For isotropic (white) noise, a mixture $p\,|\Phi^+\rangle\langle\Phi^+|+(1-p)\,I/4$, the maximum is $S_{\max}=2\sqrt2\,p$,
with no violation below $p=1/\sqrt2$, that is, above a mixed fraction $1-p=0.293$. \calc{} Decoherence in a pointer basis is not isotropic noise, so
that threshold does not describe it. If both members of the pair are massive, each branch dephases and $c=c_Ac_B$.

\paragraph{The time profile.} The rate integral of Chapter~\ref{ch:quantumrecords}, taken with its postulated rate and boundary capacity, gives an
exponential, $c(t)=e^{-t/\tau_{\rm IAM}}$. \calc{} The
alternative is the profile of the cosmic activation function, $c(t)=1-E(t/\tau_{\rm IAM})/e$ with $E(\eta)=\exp(1-1/\eta)$; that step is taken by
analogy with the cosmology, not derived. \analogy{} The two give, from Eq.~\eqref{eq:ent:smax}:
\begin{center}\small\begin{tabular}{lcccc}\toprule
$t/\tau_{\rm IAM}$ & 0.5 & 1 & 2 & 3\\\midrule
exponential: $c$, $S_{\max}$ & 0.607, 2.339 & 0.368, 2.131 & 0.135, 2.018 & 0.050, 2.002\\
assumed profile: $c$, $S_{\max}$ & 0.865, 2.644 & 0.632, 2.366 & 0.393, 2.149 & 0.283, 2.079\\\bottomrule
\end{tabular}\end{center}
\calc{} In both cases the violation fades towards 2 and is never lost at a finite time. A Bell test therefore adds nothing to a direct measurement of
$c(t)$ by interference; it loses its violation only when $c$ falls below the experiment's noise. The quantities that discriminate are the coherence
itself: its onset (zero slope for the assumed profile, finite slope for the exponential), and the scaling of $\tau_{\rm IAM}$ with temperature and mass.

\section{Where the discriminating experiment lies}\label{sec:ent:experiment}
With $\tau_{\rm IAM}=\hbar(\kB T)^2\ln2/E_G^3$ (Eq.~\eqref{eq:qr_tauIAM}, Chapter~\ref{ch:quantumrecords}), $E_G=Gm^2/R$, and the boundary capacity $\kB T/E_G$ taken as an
assumption, silica spheres ($\rho=2200$\,kg\,m$^{-3}$, fused silica) give:
\begin{center}\small\begin{tabular}{lccl}\toprule
system & $\tau_{\rm IAM}$ & $\tau_{\rm PD}=\hbar/E_G$ & reading\\\midrule
$10^{-15}$\,kg, 10\,mK & $5.1\times10^{17}$\,s & 0.75\,s & Di\'osi--Penrose testable; IAM channel not\\
$10^{-12}$\,kg, 10\,mK & 509\,s & 7.5\,$\mu$s & loss in $\mu$s, or none on that scale\\
$10^{-12}$\,kg, 20\,mK & 2{,}034\,s & 7.5\,$\mu$s & doubling $T$ quadruples $\tau_{\rm IAM}$\\
$\tau_{\rm IAM}=\tau_{\rm PD}$ at 10\,mK & \multicolumn{2}{c}{$m\approx2.2\times10^{-10}$\,kg} & crossover\\\bottomrule
\end{tabular}\end{center}
\calc{} The $10^{-12}$\,kg (nanogram) row and the crossover are those of Chapter~\ref{ch:quantumrecords}. The $10^{-15}$\,kg (picogram) row shows the other side: there the
Di\'osi--Penrose time is within reach of levitated optomechanics~\cite{RomeroIsart2011} and the IAM channel is not. The discriminating experiment is
a superposition near $10^{-12}$\,kg held at two or more bath temperatures. Di\'osi--Penrose~\cite{Diosi1989,Penrose1996} predicts loss of coherence
in microseconds at any temperature; the IAM channel predicts none on that scale, and a time that grows as $T^2$. \prediction

\section{Other entanglement results}\label{sec:ent:other}
\paragraph{Top quarks.} Entanglement of top--antitop spin states near threshold has been observed at the LHC~\cite{ATLAS2024TopEnt,CMS2024TopEnt}.
\observed{} The top quark lives $\hbar/\Gamma_t=4.6\times10^{-25}$\,s ($\Gamma_t=1.42$\,GeV~\cite{PDG2022}) and its spin is read from its decay
products. \calc{} No gravitational channel acts on that time scale. The result is consistent with the law and does not test it.

\paragraph{Violation without entanglement.} A Bell inequality has been violated by more than four standard deviations using photons from separate
sources with no entanglement, through indistinguishability by path identity in four-photon frustrated interference; the analysis uses postselected
fourfold coincidences~\cite{Wang2025SciAdv}. \observed{} This concerns the structure of the quantum state and lies outside the gravitational channel.
The law adds only the cost of each coincidence record at its detectors.

\paragraph{Photons that never coexisted.} Entanglement swapping has correlated two photons, the first detected before the second was
emitted~\cite{Megidish2013}. \observed{} Standard quantum mechanics predicts the correlation. The law writes one record at each detection, paid at that
detector, and assigns none to the interval between them.

\section{The ledger only grows}\label{sec:ent:ledger}
The cosmological side needs one more statement from the law. The accumulated record entropy is a sum of irreversible records, so it cannot fall.
Read as that ledger, $E(a)=\exp(1-1/a)$ must increase monotonically, and it does: $d\ln E/da=1/a^2>0$. \derived{} The same integral fixes how the
ledger is normalised. With the exponent $n=7/2$ of Chapter~\ref{ch:quantumrecords}, the entropy increment per unit $a$ scales as $a^{n-11/2}=a^{-2}$,
whose integral $-1/a+{\rm const}$ is dominated by its lower limit. Only the $a$-dependent part is physical, and the constant is fixed by $E(1)=1$.
\calc

\section{Status}
Stands: quantum mechanics, and with it every Bell correlation, is unchanged; the law gives each record its cost and the surface that holds it, paid
on erasure or reset at that surface's temperature; photon pairs in vacuum keep $|S|=2\sqrt2$ at any distance. \derived{} Decoherence in a pointer
basis leaves $S_{\max}=2\sqrt{1+c^2}$, which fades to 2 and is never lost at a finite time; the observable is the coherence $c(t)$. \derived{}
Predicted: $\tau_{\rm IAM}\propto T^2m^{-5}$ for massive superpositions, separable from Di\'osi--Penrose near $10^{-12}$\,kg. \prediction{} Open:
the boundary capacity $\kB T/E_G$ and the time profile of $c(t)$, which needs a derivation at the superposition boundary before it can discriminate.
\openprob
